\subsubsection{边界积分方法}\label{subsec:boundary-integral-methods}

设 $\Omega$ 为具有光滑边界的有界区域. 希望求解问题
\begin{equation}\label{eq:harmonic-dirichlet-boundary-integral-problem}
\Delta u=0,\ \bm{x}\in\Omega,\qquad u=\phi(\bm{x}),\ \bm{x}\in\partial\Omega.
\end{equation}
若已知 Green 函数, 就会有表示式
\begin{equation}\label{eq:harmonic-green-normal-boundary-representation}
u(\bm{x})=\int_{\partial\Omega}\dfrac{\partial}{\partial n_{\bm{y}}}G(\bm{x},\bm{y})\phi(\bm{y})\,\mathrm{d}S_{\bm{y}}.
\end{equation}
作与 \eqref{eq:harmonic-green-normal-boundary-representation} 类似的设定, 用 Laplace 方程的基本解代替 Green 函数
\begin{equation}\label{eq:laplace-double-layer-potential-ansatz}
u(\bm{x})=\int_{\partial\Omega}\dfrac{\partial}{\partial n_{\bm{y}}}\mathcal{G}(|\bm{x}-\bm{y}|)g(\bm{y})\,\mathrm{d}S_{\bm{y}}.
\end{equation}
这里函数 $g$ 未知, 要寻找联系 $g$ 与 $\phi$ 的方程.

注意, 对任意 $g\in C(\partial\Omega)$, 由 \eqref{eq:laplace-double-layer-potential-ansatz} 给出的函数 $u$ 在 $\Omega$ 内调和; 只需在积分号下关于 $\bm{x}$ 作 Laplace 运算即可. 为满足边界条件, 必须有
\begin{equation}\label{eq:double-layer-dirichlet-boundary-limit}
\phi(\bm{x})=\lim_{\bm{z}\to\bm{x},\bm{z}\in\Omega}\int_{\partial\Omega}\dfrac{\partial}{\partial n_{\bm{y}}}\mathcal{G}(|\bm{z}-\bm{y}|)g(\bm{y})\,\mathrm{d}S_{\bm{y}}
\end{equation}
对 $\bm{x}\in\partial\Omega$ 成立; 这就是所需的联系 $g$ 与 $\phi$ 的方程. 在 \eqref{eq:double-layer-dirichlet-boundary-limit} 中, 不能简单地以 $\bm{x}$ 替换 $\bm{z}$ 来取极限; 尽管积分在 $\bm{z}\in\partial\Omega$ 时存在, 但它在那里不连续. 事实上, 以下将证明实际有
\begin{equation}\label{eq:double-layer-interior-jump-relation}
\begin{aligned}
&\lim_{\bm{z}\to\bm{x},\bm{z}\in\Omega}\int_{\partial\Omega}\dfrac{\partial}{\partial n_{\bm{y}}}\mathcal{G}(|\bm{z}-\bm{y}|)g(\bm{y})\,\mathrm{d}S_{\bm{y}}\\
&\quad=\int_{\partial\Omega}\dfrac{\partial}{\partial n_{\bm{y}}}\mathcal{G}(|\bm{x}-\bm{y}|)g(\bm{y})\,\mathrm{d}S_{\bm{y}}+\dfrac12g(\bm{x}).
\end{aligned}
\end{equation}


回顾 Cauchy 公式也有类似情形; 若 $C$ 为平面中的光滑闭曲线且 $f$ 解析, 则
\begin{equation}\label{eq:cauchy-integral-boundary-principal-value-analogue}
\dfrac1{2\pi i}\int_C\dfrac{f(\zeta)}{\zeta-z}\,\mathrm{d}\zeta
\end{equation}
当 $z$ 在 $C$ 内部时等于 $f(z)$, 当 $z$ 在 $C$ 外部时等于 0, 而在 $C$ 上等于 $f(z)/2$ (按主值意义). 将 \eqref{eq:double-layer-interior-jump-relation} 代入 \eqref{eq:double-layer-dirichlet-boundary-limit}, 得到
\begin{equation}\label{eq:dirichlet-double-layer-boundary-integral-equation}
\phi(\bm{x})-\dfrac12g(\bm{x})=\int_{\partial\Omega}\dfrac{\partial}{\partial n_{\bm{y}}}\mathcal{G}(|\bm{x}-\bm{y}|)g(\bm{y})\,\mathrm{d}S_{\bm{y}}.
\end{equation}


这样就把偏微分方程 \eqref{eq:harmonic-dirichlet-boundary-integral-problem} 替换为等价的积分方程 \eqref{eq:dirichlet-double-layer-boundary-integral-equation}. 这有两个优点. 如将在 Chapter~\ref{chap:linear-elliptic-equations} 中看到的, 对 \eqref{eq:dirichlet-double-layer-boundary-integral-equation} 这样的积分方程建立存在性理论相当容易. 此外, 与基于 \eqref{eq:harmonic-dirichlet-boundary-integral-problem} 的数值方法相比, 基于 \eqref{eq:dirichlet-double-layer-boundary-integral-equation} 的方法处理的是空间维数较低的问题, 因而所需网格点较少. 事实上, 关于 Laplace 方程及 Stokes 方程的 ``边界元方法'' 已有大量文献.

还需验证 \eqref{eq:double-layer-interior-jump-relation}. 设 $\bm{z}$ 靠近 $\partial\Omega$, 用 $\bm{x}$ 表示 $\partial\Omega$ 上离 $\bm{z}$ 最近的点; 不失一般性, 可以选择坐标系使 $\bm{x}$ 为原点, 且 $\partial\Omega$ 的法向为第 $m$ 个坐标方向. 设 $N$ 为原点的一个邻域; 那么可将 \eqref{eq:double-layer-dirichlet-boundary-limit} 的右端分解如下:
\begin{equation}\label{eq:double-layer-local-distant-splitting}
\begin{aligned}
\int_{\partial\Omega}\dfrac{\partial}{\partial n_{\bm{y}}}\mathcal{G}(|\bm{z}-\bm{y}|)g(\bm{y})\,\mathrm{d}S_{\bm{y}}
&=\int_{\partial\Omega\cap N}\dfrac{\partial}{\partial n_{\bm{y}}}\mathcal{G}(|\bm{z}-\bm{y}|)g(\bm{y})\,\mathrm{d}S_{\bm{y}}\\
&\quad+\int_{\partial\Omega\setminus N}\dfrac{\partial}{\partial n_{\bm{y}}}\mathcal{G}(|\bm{z}-\bm{y}|)g(\bm{y})\,\mathrm{d}S_{\bm{y}}.
\end{aligned}
\end{equation}


第二项在 $\bm{z}=\bm{0}$ 处连续. 对第一项, 选择足够小的 $N$, 使 $\partial\Omega\cap N$ 可以表示为 $y_m=\phi(y_1,\ldots,y_{m-1})$; 令 $\bm{u}=(y_1,\ldots,y_{m-1})$. 这给出
\begin{equation}\label{eq:boundary-graph-normal-surface-measure}
\bm{n}_{\bm{y}}=(-\nabla\phi,1)/\sqrt{1+|\nabla\phi|^2},\quad\mathrm{d}S_{\bm{y}}=\sqrt{1+|\nabla\phi|^2}\mathrm{d}\bm{u}.
\end{equation}
可以选取 $N$, 使 $\partial\Omega\cap N=\{(\bm{u},\phi(\bm{u}))\mid|\bm{u}|<\epsilon\}$. 现在 \eqref{eq:double-layer-local-distant-splitting} 右端第一项化为
\begin{equation}\label{eq:double-layer-local-graph-kernel}
\begin{aligned}
&\int_{\{|\bm{u}|<\epsilon\}}\nabla_{\bm{y}}\mathcal{G}(|\bm{z}-(\bm{u},\phi(\bm{u}))|)\cdot(-\nabla\phi(\bm{u}),1)g(\bm{u},\phi(\bm{u}))\,\mathrm{d}\bm{u}\\
&\quad=\int_{\{|\bm{u}|<\epsilon\}}\dfrac{-\bm{u}\cdot\nabla\phi(\bm{u})+\phi(\bm{u})-z_m}
{\Omega_m(\sqrt{|\bm{u}|^2+|\phi(\bm{u})-z_m|^2})^m}g(\bm{u},\phi(\bm{u}))\,\mathrm{d}\bm{u}.
\end{aligned}
\end{equation}
注意, 若在 \eqref{eq:double-layer-local-graph-kernel} 中令 $z_m=0$, 则当 $\bm{u}\to\bm{0}$ 时, $-\bm{u}\cdot\nabla\phi(\bm{u})+\phi(\bm{u})$ 为 $|\bm{u}|^2$ 阶, 因而被积函数为 $|\bm{u}|^{-(m-2)}$ 阶, 即它可积. 尽管积分在 $\bm{z}=\bm{0}$ 时存在, 却不能在积分号下取极限 $z_m\to0$. 现在在 $z_m<0$ 的限制下考察这一极限. 需要研究的项为
\begin{equation}\label{eq:double-layer-singular-normal-kernel}
-z_m\int_{\{|\bm{u}|<\epsilon\}}\dfrac1{\Omega_m(\sqrt{|\bm{u}|^2+|\phi(\bm{u})-z_m|^2})^m}g(\bm{u},\phi(\bm{u}))\,\mathrm{d}\bm{u}.
\end{equation}


当 $|\bm{u}|$ 和 $|z_m|$ 很小时, 有
\begin{equation}\label{eq:double-layer-graph-denominator-expansion}
\dfrac1{\sqrt{|\bm{u}|^2+(\phi(\bm{u})-z_m)^2}^{\,m}}=\dfrac1{\sqrt{|\bm{u}|^2+z_m^2}^{\,m}}(1+O(|z_m|)+O(|\bm{u}|^2)),
\end{equation}
容易验证, 当 $z_m\to0$ 时, 只有首项贡献会在 \eqref{eq:double-layer-singular-normal-kernel} 中产生不连续. 因而只需考察积分
\begin{equation}\label{eq:double-layer-leading-normal-kernel}
-z_m\int_{\{|\bm{u}|<\epsilon\}}\dfrac1{\Omega_m(\sqrt{|\bm{u}|^2+z_m^2})^m}g(\bm{u},\phi(\bm{u}))\,\mathrm{d}\bm{u}.
\end{equation}


定义
\begin{equation}\label{eq:double-layer-tangential-spherical-mean}
I_r(g)=\dfrac1{r^{m-2}\Omega_{m-1}}\int_{\{|\bm{u}|=r\}}g(\bm{u},\phi(\bm{u}))\,\mathrm{d}S.
\end{equation}
在 \eqref{eq:double-layer-leading-normal-kernel} 中作代换 $\bm{u}=-z_m\bm{v}$, 得到表达式
\begin{equation}\label{eq:double-layer-scaled-radial-integral}
\begin{aligned}
&\int_{\{|\bm{v}|\leq-\epsilon/z_m\}}\dfrac1{\Omega_m(\sqrt{|\bm{v}|^2+1})^m}g(-z_m\bm{v},\phi(-z_m\bm{v}))\,\mathrm{d}\bm{v}\\
&\quad=\dfrac{\Omega_{m-1}}{\Omega_m}\int_0^{-\epsilon/z_m}\dfrac{r^{m-2}}{(r^2+1)^{m/2}}I_{-z_mr}(g)\,\mathrm{d}r.
\end{aligned}
\end{equation}
在极限 $z_m\to0-$ 下, 得到
\begin{equation}\label{eq:double-layer-half-density-jump-limit}
\dfrac{\Omega_{m-1}}{\Omega_m}g(\bm{0})\int_0^\infty\dfrac{r^{m-2}}{(r^2+1)^{m/2}}\,\mathrm{d}r=\dfrac12g(\bm{0}).
\end{equation}


这里用到了
\begin{equation}\label{eq:double-layer-radial-gamma-integral}
\int_0^\infty\dfrac{r^{m-2}}{(r^2+1)^{m/2}}\,\mathrm{d}r=\dfrac{\Gamma(\frac{m-1}2)\sqrt\pi}{2\Gamma(\frac m2)}
\end{equation}
(见 \cite[p.~292]{GR}), 以及在 Problem~\ref{ex:unit-ball-volume-gamma-function} 中得到的 $\Omega_m$ 的表达式.
